\scanpage{133}
\mspage{78}
% archive sheet number: 132

donc il n'y a qu'une sorte de \inserted{\unsure{cat|cas}} \unsure{localisées|localisés} à examiner. On trouve ainsi

\diagram{\begin{tikzcd}[ampersand replacement=\&,column sep=large]
(\underline{\mathrm{Fibsc}}_0\,S)\deleted{\unsure{\Sigma}}W_S^{-1}\arrow[r,"\approx"]\&{\unsure{\mathcal{F}|F}W_S^{-1}}\arrow[l,bend left=16,"\approx"',"{\Phi_{0\mathcal F}}"]
\end{tikzcd}}

\diagram{\begin{tikzcd}[ampersand replacement=\&,column sep=large]
\mathcal{F}'\deleted{\unsure{\Sigma}}W_S^{-1}\arrow[r,"\approx"]\&\mathcal{F}\deleted{\unsure{\Sigma}}W_S^{-1}\arrow[l,bend left=16,"\approx"',"{\Phi_{0\mathcal{F}',\mathcal{F}}}"]
\end{tikzcd}}

pour $\mathcal{F}$ (ou $\mathcal{F},\mathcal{F}'$) sous-catégorie(s) pleine(s) de $\deleted{\unsure{\mathcal{F}}}\underline{\mathrm{Fibhot}}\,S$ telle que

\[
\underline{\mathrm{Fib}}_0\,S\subset\mathcal{F}\subset\underline{\mathrm{Fibhot}}\,S
\qquad\text{resp.}\qquad
\underline{\mathrm{Fib}}_0\,S\subset\mathcal{F}'\subset\mathcal{F}\subset\underline{\mathrm{Fibhot}}\,S,
\]

et de même du côté droit pour les sous-catégories \inserted{pl.} $\mathcal{F}$ ou $\mathcal{F},\mathcal{F}'$ de $\underline{\mathrm{Fibhot}}\,S$ satisfaisant

\[
\underline{\mathrm{Cofib}}_0\,S\subset\mathcal F\subset\underline{\mathrm{Fibhot}}\,S
\qquad\text{respt.}\qquad
\underline{\mathrm{Cofib}}_0\,S\subset\mathcal F'\subset\mathcal F\subset\underline{\mathrm{Fibhot}}\,S.
\]

\inserted{\uline{NB} La catégorie $\underline{\mathrm{Parf}}\,S$ contient à la fois $\underline{\mathrm{Fib}}_0\,S$ et $\underline{\mathrm{Cofib}}_0\,S$, et fait donc partie des catégories admissibles \unsure{dont varie|dont l'une est} $\mathcal{F}$, ou $\mathcal{F}'$.}

Ainsi, les quatre foncteurs $\Phi,\Psi,\Phi_0,\Psi_0$ donnent tous les quasi-inverses qu'on peut souhaiter au \deleted{\unsure{premier étage}} \inserted{rez de chaussée} d'une part ($\Phi$ et $\Psi$), à l'étage de l'autre ($\Phi_0,\Psi_0$) \deleted{\unsure{(ces catégories, cf. ci-dessus)}}. Ils ramènent \unsure{tous les localisés pertinents|toutes les localisées pertinentes} aux \deleted{\unsure{trois}} \inserted{quatre} suivants

\[
\begin{aligned}
&(\mathrm{Cat}_{/S})W_S^{-1},(\mathrm{Cat}_{/S})(W_S^d)^{-1},(\mathrm{Cat}_{/S})(W_S^g)^{-1}\\
&\hspace{11em}\text{(au rez de chaussée)}\\
&(\underline{\mathrm{Fibhot}}\,S)W_S^{-1}\qquad\text{(à l'étage),}
\end{aligned}
\]

\textcircled{VII} \deleted{\unsure{après un choix|ap. un choix} \unsure{principe|privilégié} \unsure{trouvés|trois} \unsure{pleins|trois} à l'étage,} Il reste à comparer le localisé à l'étage avec les localisés \deleted{\unsure{grossiers}} au rez de chaussée. Les flèches verticales \deleted{\unsure{gauches}} descendant de l'étage au rez de chaussée sont compatibles avec $W_S^d$ du côté gauche du diagramme, avec $W_S^g$ du côté droit, donc elles induisent \inserted{essentiellement} \uline{deux} foncteurs différents
